% BibTeX keys used in this section:
%   Makoviychuk2021IsaacGym, Narang2022Factory, Tang2023IndustReal,
%   Allshire2022TriFinger, Mittal2023Orbit, Mittal2025IsaacLab,
%   SerranoMunoz2023skrl, Handa2023DeXtreme,
%   Shahid2023CurriculumTendon, Dhakate2025CaRoSaC, Caggiano2022MyoSuite,
%   Saito2022ContinuumRL, Lilge2021BenchmarkTDCR, Garrido2024CDPR,
%   Yoshikawa1985Manipulability, Rastegar1990WorkspaceMonteCarlo,
%   Peidro2017GaussianGrowth, Vahrenkamp2015ConstrainedManipulability,
%   Cao2017WorkspaceTDCR,
%   Klein2023, Nemoto2022CableWrist, Walter2023Tensegrity, Walter2022CDPRCPS,
%   Schulman2017PPO, SuttonBarto2018, Todorov2012MuJoCo

%===================================================================%
\section{Prior Work}\label{sec:prior_work}
%===================================================================%

This section reviews prior work along the three dimensions that converge in
the present thesis: (i)~\acl{RL}-based manipulation in \ac{GPU}-accelerated
simulators, (ii)~simulation and control of tendon- and cable-driven robots,
and (iii)~workspace analysis methods for custom manipulators.
A final subsection summarizes FAPS-internal predecessor work.
The review focuses on contributions that directly inform the methodological
choices made in this thesis and concludes by identifying the specific gap that
the present work addresses.

%-------------------------------------------------------------------
\subsection{GPU-accelerated reinforcement learning for manipulation}
\label{subsec:prior_gpu_rl}
%-------------------------------------------------------------------

\paragraph{Isaac Gym and large-scale parallel training.}
Makoviychuk~\emph{et~al.}\ demonstrated that running both PhysX simulation and
neural-network training entirely on GPU eliminates CPU--GPU data transfer
bottlenecks and yields two to three orders of magnitude training speedup
over CPU-based simulators~\cite{Makoviychuk2021IsaacGym}.
In their Franka cube stacking benchmark, 16\,384 parallel agents trained with
\ac{PPO} converged in under 25~minutes on a single A100 GPU, and a Shadow Hand
cube rotation task achieved 20 consecutive successful rotations in
approximately one hour~\cite{Makoviychuk2021IsaacGym}.
These results established Isaac Gym as the foundational platform for
GPU-accelerated \ac{RL} and underpin the Isaac Sim/IsaacLab ecosystem
used in the present work.

\paragraph{Contact-rich manipulation and assembly.}
Narang~\emph{et~al.}\ introduced \emph{Factory}, a suite of Isaac Gym
environments for contact-rich assembly tasks that achieves a reported
20\,000-fold speedup over prior state-of-the-art contact simulation
through signed-distance-function collisions, contact reduction from
approximately one million naive contacts to approximately 10\,000, and a
Gauss--Seidel solver~\cite{Narang2022Factory}.
Building on Factory, Tang~\emph{et~al.}\ (\emph{IndustReal}) transferred
peg-insertion and gear-assembly policies from simulation to a real Franka
Panda, reporting 83--99\% success rates over 600 physical trials with
realistic part clearances of $\leq 0.5$--$0.6\,\text{mm}$~\cite{Tang2023IndustReal}.
IndustReal introduced a simulation-aware policy update mechanism that detects
unreliable simulator states and a sampling-based curriculum for progressive
task difficulty---techniques that are directly adoptable for conveyor-based
pick-and-place policy design~\cite{Tang2023IndustReal}.

\paragraph{Sim-to-real transfer for dexterous manipulation.}
Allshire~\emph{et~al.}\ achieved an 83\% real-world success rate for
6-\ac{DoF} cube reposing on a TriFinger robot, trained entirely in simulation
on a single GPU using \ac{PPO} with asymmetric actor-critic and
keypoint-based pose representations~\cite{Allshire2022TriFinger}.
The keypoint representation was shown to provide more stable gradients than
position-plus-quaternion encoding, and domain randomization over friction,
mass, and sensor noise enabled zero-shot sim-to-real transfer with 16\,384
parallel environments~\cite{Allshire2022TriFinger}.
Handa~\emph{et~al.}\ (\emph{DeXtreme}) extended this line of work to agile
in-hand manipulation with an Allegro Hand using automatic domain
randomization, which dynamically adjusts randomization ranges without manual
tuning~\cite{Handa2023DeXtreme}.
Both works employ \ac{PPO} with asymmetric actor-critic, providing the
privileged critic with full simulation state while limiting the actor to
real-world-available observations~\cite{Allshire2022TriFinger,Handa2023DeXtreme}.

\paragraph{Frameworks: Orbit, Isaac Lab, and skrl.}
Mittal~\emph{et~al.}\ presented \emph{Orbit}, a modular design powered by
NVIDIA Isaac Sim with GPU-based parallelization, wrapping four \ac{RL}
libraries (rl\_games, RSL-RL, Stable-Baselines3, skrl) and providing
20+ benchmark tasks~\cite{Mittal2023Orbit}.
Orbit evolved into \emph{Isaac Lab} in March~2024, achieving up to
1.6~million frames per second for state-based manipulation tasks and
adding mimic joint support, closed-loop kinematic chains, and a
deformable-object solver~\cite{Mittal2025IsaacLab}.
The \texttt{skrl} library provides native first-class support for Isaac Lab
environments with competitive performance to rl\_games on standard
benchmarks while offering greater modularity through seven standalone
components~\cite{SerranoMunoz2023skrl}.
In the present thesis, Isaac Lab serves as the development framework and
skrl provides the \ac{PPO} training backend.

%-------------------------------------------------------------------
\subsection{Tendon- and cable-driven robot simulation}
\label{subsec:prior_tendon_sim}
%-------------------------------------------------------------------

\paragraph{Tendon-driven RL in MuJoCo.}
Shahid~\emph{et~al.}\ modeled a bio-inspired ``Robostrich'' arm using
MuJoCo's native tendon elements with attachment sites routed through
predefined wire run-through holes, demonstrating that general actuators
acting on tendons produce realistic tendon-contraction-to-joint-actuation
coupling~\cite{Shahid2023CurriculumTendon}.
The study showed that curriculum-based \ac{RL} outperforms vanilla SAC for
reaching tasks on this underactuated tendon-driven manipulator by
progressing from simple to complex targets~\cite{Shahid2023CurriculumTendon}.
Saito and Morimoto presented the first systematic \ac{RL} comparison between
a 9-tendon continuum arm and a 7-\ac{DoF} rigid manipulator across multiple
tasks (reaching, crank rotation, throwing, peg-in-hole), finding that
continuum arms excel in uncertainty-rich and anisotropic tasks while rigid
arms are superior for precision tasks~\cite{Saito2022ContinuumRL}.
Sim-to-real validation on a real pneumatic continuum robot confirmed the
MuJoCo-based model's validity~\cite{Saito2022ContinuumRL}.

\paragraph{Musculoskeletal simulation.}
Caggiano~\emph{et~al.}\ (\emph{MyoSuite}) provided an automated pipeline
converting OpenSim musculoskeletal models to MuJoCo, validating forward
kinematics, moment arms, and muscle forces across 39 muscle-tendon units
in MyoHand (23 joints, 29 bones)~\cite{Caggiano2022MyoSuite}.
MyoSuite achieves over two orders of magnitude speedup compared to OpenSim
and offers 204 tasks spanning dexterous manipulation with antagonistic
muscle-tendon actuation, including non-stationarity scenarios such as
muscle fatigue and tendon-transfer surgery~\cite{Caggiano2022MyoSuite}.
The OpenSim-to-MuJoCo validation methodology is relevant as a reference
for validating the Isaac Sim tendon model developed in this thesis.

\paragraph{Cable-driven parallel robot simulation with RL.}
Dhakate~\emph{et~al.}\ (\emph{CaRoSaC}) used Unity3D with Obi Rope for
flexible cable simulation via Extended Position-Based Dynamics, validated
against a real industrial 4-cable-suspended parallel robot~\cite{Dhakate2025CaRoSaC}.
The TD3-based \ac{RL} controller outperformed classical kinematic solvers
in trajectory tracking, especially near workspace boundaries where
force-distribution assumptions break down~\cite{Dhakate2025CaRoSaC}.
A two-stage training approach---first in a simplified no-sag environment,
then in a full simulation with realistic cable sag---reduced training time
while maintaining performance~\cite{Dhakate2025CaRoSaC}.
Garrido~Campos~\emph{et~al.}\ validated a cable-driven parallel robot
simulation against a hardware prototype, achieving a peak position
divergence of only 0.27\% across various speed
scenarios~\cite{Garrido2024CDPR}.

\paragraph{Tendon-driven modeling benchmarks.}
Lilge and Burgner-Kahrs comprehensively benchmarked four tendon-driven
continuum robot modeling approaches---constant curvature, pseudo-rigid-body,
and Cosserat rod formulations---under varying loads~\cite{Lilge2021BenchmarkTDCR}.
A key finding is that tendon force expressions differ across the literature
and yield divergent results under external loading, establishing
accuracy-versus-computation-time trade-offs that inform which level of
fidelity a physics engine should reproduce~\cite{Lilge2021BenchmarkTDCR}.

%-------------------------------------------------------------------
\subsection{Workspace analysis methods}
\label{subsec:prior_workspace}
%-------------------------------------------------------------------

\paragraph{Manipulability measures.}
Yoshikawa defined manipulability as $w = \sqrt{\det(\mathbf{J}\mathbf{J}^\top)}$,
where $\mathbf{J}$ is the manipulator Jacobian, quantifying how well a
manipulator can change end-effector position and orientation from a given
configuration~\cite{Yoshikawa1985Manipulability}.
The resulting manipulability ellipsoid has the singular values of
$\mathbf{J}$ as semi-axes, and $w = 0$ identifies singular
configurations~\cite{Yoshikawa1985Manipulability}.
Vahrenkamp and Asfour extended this measure by penalizing Jacobian entries
based on proximity to joint limits and self-collision distances, yielding
a constrained manipulability metric that naturally reduces near workspace
boundaries~\cite{Vahrenkamp2015ConstrainedManipulability}.
This constraint-penalized approach is directly applicable to tendon-driven
manipulators, where cable tension limits act analogously to joint
limits~\cite{Vahrenkamp2015ConstrainedManipulability}.

\paragraph{Computational workspace determination.}
Rastegar and Fardanesh established the Monte Carlo forward-kinematics
sampling method for workspace determination, in which random joint-angle
vectors are uniformly sampled within joint limits and mapped through
forward kinematics to produce a workspace point
cloud~\cite{Rastegar1990WorkspaceMonteCarlo}.
A known limitation is non-uniform density: uniform joint-space sampling
produces denser coverage near the workspace center and sparser coverage
at boundaries~\cite{Rastegar1990WorkspaceMonteCarlo}.
Peidr\'{o}~\emph{et~al.}\ addressed this limitation with a Gaussian Growth
method that iteratively densifies boundary regions, achieving significantly
more accurate workspace boundaries for high-\ac{DoF} robots at comparable
computation cost~\cite{Peidro2017GaussianGrowth}.
Cao~\emph{et~al.}\ performed workspace analysis specifically on
tendon-driven continuum robots, determining maximum bending angles from
physical interference between tendon paths and the elastic backbone, and
using a superposition method to combine individual module configuration
spaces~\cite{Cao2017WorkspaceTDCR}.

%-------------------------------------------------------------------
\subsection{Identified gap}
\label{subsec:prior_gap}
%-------------------------------------------------------------------

The literature reviewed above reveals a clear separation between two bodies
of work.
On one side, GPU-accelerated \ac{RL} for manipulation has advanced rapidly
within the Isaac Gym / Isaac Sim ecosystem, with demonstrated sim-to-real
transfer for contact-rich tasks involving rigid-body
robots~\cite{Makoviychuk2021IsaacGym,Narang2022Factory,Tang2023IndustReal,Allshire2022TriFinger,Handa2023DeXtreme}.
On the other side, tendon- and cable-driven robot simulation for \ac{RL}
has been explored predominantly in
MuJoCo~\cite{Shahid2023CurriculumTendon,Caggiano2022MyoSuite,Saito2022ContinuumRL,Todorov2012MuJoCo}
or in specialized engines~\cite{Dhakate2025CaRoSaC}.
No published work, as of March~2026, combines tendon-driven robot actuation
modeling with \ac{RL}-based policy training in NVIDIA Isaac Sim or Isaac Lab.
Furthermore, the FAPS-internal predecessor work by Klein~\cite{Klein2023}
established workspace analysis and simulation for a cable-driven manipulator
but did not employ learning-based control.
The present thesis addresses this gap by constructing a validated tendon-driven
robot model in Isaac Sim, formulating progressive \ac{RL} tasks within
Isaac Lab, and training manipulation policies using
\texttt{skrl}~\ac{PPO}~\cite{SerranoMunoz2023skrl,Schulman2017PPO}.
